\section{Conclusions}
\label{sec:conclusion}

\begin{enumerate}
\item HEAs with FCC CoCrFeMnNi-type or Mo-modified eutectic compositions exhibit passive-film corrosion resistance in neutral chloride media that matches or exceeds the 316L benchmark, making them credible candidates for seawater-wetted marine components \cite{qiu2017,lpbf_cantor_2024,mo_ehea_2026}.
\item Additive manufacturing is not merely a shaping tool but a \emph{corrosion-engineering} tool for HEAs: rapid solidification suppresses elemental segregation, refines microstructure, and measurably improves pitting resistance over cast material---provided defects and post-build segregation are controlled \cite{slm_ehea_2025,han2025,cagirici2024}.
\item Process selection maps onto marine needs: LPBF for complex small components in seawater systems, SEBM for crack-prone and refractory systems, DED for cladding, repair, and functionally graded construction, and WAAM for large structures once HEA wire feedstock matures \cite{cagirici2024,raultaiwade2021,jiang2019clad}.
\item Corrosion performance is environment-specific: AM CoCrFeMnNi surpasses 316L in chloride but is far inferior in acidic media, so crevice, under-deposit, and biofilm microenvironments must be tested explicitly rather than extrapolated from NaCl screening \cite{lpbf_cantor_2024}.
\item The decisive gaps for deployment are evidential and regulatory: long-term natural seawater exposure, corrosion fatigue/SCC, biofilm coupling, and galvanic/CP compatibility data are scarce, and no marine AM standard yet accommodates HEA compositions \cite{dnvstb203,abs_mam,lr_am}.
\item The practical near-term pathway is hybrid: HEA claddings and graded deposits on marine steel, qualified through a staged ladder of electrochemical screening, long-term exposure, coupled-degradation testing, and system-level galvanic assessment.
\end{enumerate}
